\documentclass[11pt,a4paper]{article}

\usepackage[utf8]{inputenc}
\usepackage[T1]{fontenc}
\usepackage[spanish,es-tabla]{babel}

\usepackage{helvet} 
\renewcommand{\familydefault}{\sfdefault}

\usepackage[left=2.5cm, right=2.5cm, top=1.5cm, bottom=1.5cm]{geometry}
\usepackage{parskip}

\usepackage{array}
\usepackage{xcolor}
\usepackage{enumitem}
\usepackage{graphicx}

\usepackage{fancyhdr}
\usepackage{lastpage}
\pagestyle{fancy}
\fancyhf{}
\renewcommand{\headrulewidth}{0.4pt}
\renewcommand{\footrulewidth}{0pt}

\lhead{\textcolor{gray}{\small Título de mi documento}} 
\rhead{\textcolor{gray}{\small \today}} 
\cfoot{\small Página \thepage\ de \pageref*{LastPage}} 

\usepackage[colorlinks=true, linkcolor=blue, urlcolor=blue]{hyperref}


\begin{document}

\section*{Introducción}
Esta es tu plantilla mejorada. Gracias al paquete \texttt{parskip}, los párrafos ya no tienen sangría, sino que se separan por un espacio vertical, dándole un aspecto mucho más limpio.

\section*{Listas y Estilos}
Al tener \texttt{enumitem} y \texttt{xcolor}, puedes hacer listas estilizadas rápidamente:
\begin{itemize}[leftmargin=*]
    \item Primer elemento de la lista.
    \item Segundo elemento con un \href{https://es.wikipedia.org}{enlace azul interactivo}.
\end{itemize}

\end{document}